Spatiotemporal graph neural networks (STGNNs) dominate traffic forecasting benchmarks,
yet which architectural properties govern their robustness under realistic deployment
shifts remains unclear.
Eight frozen architectures are evaluated under three deployment-motivated shifts---inference-time
sensor dropout, zero-shot cross-city transfer (METR-LA\,$\to$\,PEMS-BAY), and
weekday/weekend temporal subgroup shift---with all predictions pre-registered before
any experiments ran.

Robustness turns out to be shift-specific, and architectural progress changes failure
modes rather than eliminating them.
Adjacency routing concentration, $\mathrm{max\_mass}(A)$, explains the sensor-dropout
ranking at Pearson $\rho{=}{+}0.979$ across all four adaptive models.
Graph WaveNet's dense softmax routing ($\mathrm{max\_mass}{\approx}0.23$) degrades
$4.3{\times}$ more than DCRNN under 30\% dropout; a mask-aware inference variant that
zeroes dropped-sensor adjacency entries recovers 75\% of the gap without retraining.
Sensor-dropout robustness is therefore not a property of topology learning per se, but of
routing density: AGCRN ($\mathrm{max\_mass}{=}0.008$) and STAEformer show fragility
comparable to DCRNN despite learning adaptive spatial relations.
Cross-city and temporal fragility lie on orthogonal axes: MegaCRN achieves the best
zero-shot city-transfer score yet exhibits $8$--$9{\times}$ larger temporal subgroup gaps
than the canonical bank---a fragility axis invisible in five-model evaluations and not
driven by explicit temporal covariates in MegaCRN (mechanism ablations,
\S\ref{sec:mechanism}).
Deployment-oriented model selection should be based on the target shift type and routing
density of candidate architectures, not on in-distribution accuracy alone.
